\chapter{IMPLEMENTATION}
\label{ch:implementation}

This chapter presents the software implementation of the
\texttt{kfc\_procedure} library. The repository implements two connected machine learning subsystems: the KFC clusterwise learning framework and the COBRA prediction-space aggregation framework. These subsystems are implemented as a single Python package so that clusterwise local learning and kernel-based aggregation can be used independently or combined within one workflow.

Building on the methodological design introduced in Chapter~\ref{ch:methodology}, this chapter focuses on how the framework is realized in code. The discussion covers package organization,
interface design, factory-based component registration, clustering and
divergence implementations, COBRA aggregation components, orchestration
logic, utility functions, tests, and packaging infrastructure.

The implementation follows a research-library structure rather than a
large application architecture. It uses a \texttt{src/} package layout,
scikit-learn-style estimator interfaces, modular factories, and
component-level tests for the COBRA core subsystem. The chapter therefore
emphasizes traceability between the repository structure and the
implemented machine learning workflows.

% =============================================================
\section{Repository Architecture}
\label{sec:repository_architecture}

The repository is organized as a Python package using a \texttt{src/}
layout. The main implementation is located under
\texttt{src/kfc\_procedure/}, while experiment notebooks, tests, package
configuration, and CI workflow files are stored at the repository level.

\subsection{Top-Level Repository Structure}

The main repository contains the following verified components:

\begin{verbatim}
kfc-procedure/
├── .github/workflows/ci-cd.yml     # this auto
├── notebooks/
│   ├── kfc.ipynb
│   └── experiments/
│       ├── classification/
│       └── regression/
├── src/
│   └── kfc_procedure/
├── tests/
│   └── cobra/
├── pyproject.toml
├── requirements.txt
├── README.md
├── LICENSE
└── .pylintrc
\end{verbatim}

The \texttt{pyproject.toml} file defines the package metadata, Python
requirement, dependencies, optional extras, setuptools configuration,
and pytest settings. The package name is \texttt{kfc-procedure}, with
the import package located at \texttt{kfc\_procedure}. The project uses
Python~3.11 or newer.

\subsection{Main Package Structure}

The implementation package is organized as follows:

\begin{verbatim}
kfc_procedure/
├── __init__.py
├── kfc.py
├── core/
│   ├── clustering/
│   ├── combiner/
│   ├── ml/
│   ├── steps/
│   └── factory.py
├── cobra/
│   ├── gradientcobra.py
│   ├── mixcobra.py
│   ├── combined_classifier.py
│   ├── superlearner.py
│   ├── core/
│   └── utils/
└── utils/
    ├── logger.py
    └── resolve.py
\end{verbatim}

This structure separates the framework into four implementation layers:

\begin{itemize}
      \item \textbf{KFC core layer:} Implemented under
            \texttt{core/}. It contains Bregman clustering, local model
            wrappers, K-step/F-step/C-step workflow components, and KFC
            combiners.

      \item \textbf{COBRA aggregation layer:} Implemented under
            \texttt{cobra/}. It contains \texttt{GradientCOBRA},
            \texttt{MixCOBRARegressor}, \texttt{CombinedClassifier},
            \texttt{SuperLearner}, and reusable COBRA core components.

      \item \textbf{Orchestration layer:} Implemented in
            \texttt{kfc.py}. It exposes the high-level classes
            \texttt{KFCProcedure}, \texttt{KFCRegressor}, and
            \texttt{KFCClassifier}.

      \item \textbf{Utility layer:} Implemented under both
            \texttt{kfc\_procedure/utils/} and \texttt{kfc\_procedure/cobra/utils/}.
            The top-level utilities support KFC logging and resolution helpers,
            while COBRA utilities support splitting, normalization constants,
            estimator resolution, and distance helpers.
\end{itemize}

\subsection{Clusterwise Learning Subsystem}

The KFC subsystem is implemented under \texttt{core/}:

\begin{verbatim}
core/
├── clustering/
│   ├── bregman.py
│   └── divergences/
├── ml/
├── steps/
│   ├── kstep.py
│   ├── fstep.py
│   └── cstep.py
├── combiner/
└── factory.py
\end{verbatim}

The \texttt{clustering/} package implements \texttt{BregmanKMeans} and
the divergence classes used for cluster assignment. The \texttt{steps/}
package implements the three KFC workflow stages. The \texttt{ml/}
package wraps scikit-learn estimators for local cluster-level learning,
and the \texttt{combiner/} package implements final aggregation
strategies.

\subsection{COBRA Aggregation Subsystem}

The COBRA subsystem is implemented under \texttt{cobra/}:

\begin{verbatim}
cobra/
├── gradientcobra.py
├── mixcobra.py
├── combined_classifier.py
├── superlearner.py
├── core/
│   ├── adapters/
│   ├── aggregators/
│   ├── cv/
│   ├── distances/
│   ├── estimators/
│   ├── kernels/
│   ├── losses/
│   ├── normalizers/
│   ├── optimizers/
│   ├── splitters/
│   ├── factory.py
│   └── types.py
└── utils/
\end{verbatim}

The main COBRA classes coordinate reusable core components such as
distances, kernels, losses, optimizers, and aggregators. This design
allows \texttt{GradientCOBRA}, \texttt{MixCOBRARegressor}, and
\texttt{CombinedClassifier} to share the same component infrastructure
while implementing different aggregation workflows.

\subsection{Architectural Summary}

The repository implements a two-subsystem architecture. The KFC subsystem
builds divergence-specific clusterwise predictors, while the COBRA
subsystem aggregates predictions in prediction space using
distance-based and kernel-based mechanisms. The two subsystems interact
through the KFC C-step, where COBRA-based combiners can be selected as
aggregation strategies.


\section{Software Design Principles}
\label{sec:software_design_principles}

The \texttt{kfc\_procedure} framework is designed based on a set of
established software engineering principles that ensure modularity,
extensibility, maintainability, and interoperability. Given the
complexity of combining clusterwise learning with prediction-space
aggregation, the system architecture is carefully structured to enforce
clear separation of concerns and well-defined abstraction boundaries.

This section presents the key design principles used throughout the
implementation, together with illustrative diagrams that describe their
application in the system.

% ============================================================
\subsection{Layered and Modular Architecture}

The framework follows a layered architecture in which high-level
orchestration is separated from low-level algorithmic components.

\begin{figure}[H]
      \centering
      \begin{tikzpicture}[
                  node distance=1.4cm,
                  box/.style={
                              draw,
                              rounded corners,
                              align=center,
                              text width=4.5cm,
                              font=\small
                        },
                  arrow/.style={-Latex, thick}
            ]

            \node[box] (top) {Orchestration Layer\\ \texttt{kfc.py}};

            \node[box, below left=of top] (kfc) {Clusterwise Learning\\ \texttt{core/}};
            \node[box, below right=of top] (cobra) {COBRA Aggregation\\ \texttt{cobra/}};

            \node[box, below=of kfc] (kfc2) {
                  Clustering,\\ ML Models, \\and Pipeline Steps
            };

            \node[box, below=of cobra] (cobra2) {
                  Kernels,\\ Optimizers, \\Distances, Aggregators
            };

            \draw[arrow] (top) -- (kfc);
            \draw[arrow] (top) -- (cobra);

            \draw[arrow] (kfc) -- (kfc2);
            \draw[arrow] (cobra) -- (cobra2);

      \end{tikzpicture}
      \caption{Layered modular architecture of the \texttt{kfc\_procedure} framework}
\end{figure}

\textbf{Explanation.}
This structure ensures that the orchestration layer coordinates the
pipeline, while the underlying learning and aggregation components
operate independently. This reduces coupling and improves scalability.

% ============================================================
\subsection{Interface-Driven Design (Abstraction)}

All major components are defined using abstract base classes, ensuring
that implementations follow a consistent contract.

\begin{figure}[H]
      \centering
      \begin{tikzpicture}[
                  node distance=1.6cm,
                  box/.style={draw, rounded corners, align=center, minimum width=3.8cm},
                  arrow/.style={-Latex, thick}
            ]

            \node[box] (base) {BaseEstimator};

            \node[box, below left=of base] (a) {Sklearn Wrapper};
            \node[box, below=of base] (b) {Mean Regressor};
            \node[box, below right=of base] (c) {Custom Model};

            \draw[arrow] (base) -- (a);
            \draw[arrow] (base) -- (b);
            \draw[arrow] (base) -- (c);

      \end{tikzpicture}
      \caption{Interface-driven design using abstract base estimators}
\end{figure}

\textbf{Explanation.}
This design allows multiple implementations to be swapped without
modifying the overall pipeline, ensuring flexibility and consistency
across modules such as estimators, kernels, and optimizers.

% ============================================================
\subsection{Composition over Inheritance}

Instead of deep inheritance hierarchies, the system relies on
composition to build complex pipelines from independent components.

\begin{figure}[H]
      \centering
      \begin{tikzpicture}[
                  box/.style={draw, rounded corners, align=center, minimum width=4.2cm},
                  arrow/.style={-Latex, thick}
            ]

            % Main system
            \node[box] (kfc) {KFCProcedure};

            % Core pipeline steps
            \node[box, below left=of kfc] (kstep) {KStep\\Divergence-based Clustering};
            \node[box, below=of kfc] (fstep) {FStep\\Local Models per Cluster};
            \node[box, below right=of kfc] (cstep) {CStep\\Combiner / Aggregation};

            % Internal components of each step
            \node[box, below=1.4cm of kstep] (div) {Bregman Divergences\\(Multiple Metrics)};
            \node[box, below=1.4cm of fstep] (est) {BaseEstimator\\(Local Learners)};
            \node[box, below=1.4cm of cstep] (comb) {Combiner Strategy\\(Voting / Weighted / Stacking)};

            % Inheritance layer
            \node[box, left=1.4cm of kfc] (sk) {BaseEstimator\\(scikit-learn API)};

            % Connections (composition)
            \draw[arrow] (kfc) -- (kstep);
            \draw[arrow] (kfc) -- (fstep);
            \draw[arrow] (kfc) -- (cstep);

            % Internal decomposition
            \draw[arrow] (kstep) -- (div);
            \draw[arrow] (fstep) -- (est);
            \draw[arrow] (cstep) -- (comb);

            % Inheritance relation (different meaning)
            \draw[arrow] (kfc) -- (sk);

      \end{tikzpicture}
      \caption{Hybrid composition-based pipeline with scikit-learn inheritance interface.}
\end{figure}

\textbf{Explanation.}
Each component is injected into the main pipeline, allowing independent
replacement and reducing tight coupling between modules.

% ============================================================
\subsection{Strategy Design Pattern}

Different algorithms are implemented as interchangeable strategies.

\begin{figure}[H]
      \centering
      \begin{tikzpicture}[
                  box/.style={draw, rounded corners, align=center, minimum width=3.6cm},
                  arrow/.style={-Latex, thick}
            ]

            \node[box] (s) {Aggregation Strategy};

            \node[box, below left=of s] (m) {Weighted Mean};
            \node[box, below=of s] (v) {Weighted Vote};
            \node[box, below right=of s] (g) {GradientCOBRA};

            \draw[arrow] (s) -- (m);
            \draw[arrow] (s) -- (v);
            \draw[arrow] (s) -- (g);

      \end{tikzpicture}
      \caption{Strategy pattern for aggregation methods}
\end{figure}

\textbf{Explanation.}
This design allows runtime selection of different aggregation methods
without modifying the core pipeline logic.

% ============================================================
\subsection{Factory-Based Construction}

Components are dynamically created using a factory layer.

\begin{figure}[H]
      \centering
      \begin{tikzpicture}[
                  box/.style={draw, rounded corners, align=center, minimum width=4cm},
                  arrow/.style={-Latex, thick}
            ]

            \node[box] (cfg) {Configuration Input};

            \node[box, below=of cfg] (f) {Factory Layer\\ \texttt{factory.py}};

            \node[box, below left=of f] (c1) {Clusterer};
            \node[box, below=of f] (c2) {Estimator};
            \node[box, below right=of f] (c3) {Kernel / Optimizer};

            \draw[arrow] (cfg) -- (f);
            \draw[arrow] (f) -- (c1);
            \draw[arrow] (f) -- (c2);
            \draw[arrow] (f) -- (c3);

      \end{tikzpicture}
      \caption{Factory-based dynamic component creation}
\end{figure}

\textbf{Explanation.}
This enables configuration-driven model construction and improves
flexibility in experimental setups.

% ============================================================
\subsection{Separation of Concerns}

Each module is responsible for a single well-defined functionality.

\begin{figure}[H]
      \centering
      \begin{tikzpicture}[
                  box/.style={
                        draw,
                        rounded corners,
                        align=center,
                        text width=3.2cm,
                        font=\small,
                        minimum height=1.4cm
                  },
                  center/.style={
                        draw,
                        rounded corners,
                        align=center,
                        text width=4.5cm,
                        fill=gray!10,
                        font=\small,
                        minimum height=1.5cm
                  },
                  arrow/.style={-Latex, thick}
            ]

            % Center system
            \node[center] (sys) {KFCProcedure System};

            % Independent modules
            \node[box, above=1.5cm of sys] (a) {core/clustering};
            \node[box, left=1.6cm of sys] (b) {core/ml};
            \node[box, right=1.6cm of sys] (c) {cobra/kernels};
            \node[box, below left=1.4cm of sys] (d) {cobra/optimizers};
            \node[box, below right=1.4cm of sys] (e) {utils};

            % connections (usage, not flow)
            \draw[arrow] (sys) -- (a);
            \draw[arrow] (sys) -- (b);
            \draw[arrow] (sys) -- (c);
            \draw[arrow] (sys) -- (d);
            \draw[arrow] (sys) -- (e);

      \end{tikzpicture}
      \caption{Modular separation of concerns: independent components used by the KFCProcedure system}
\end{figure}

\textbf{Explanation.}
This modular separation enforces clear responsibility boundaries across subsystems, enabling independent development, testing, and extension of each component.

% ============================================================
\subsection{Extensibility-Oriented Design}

The KFCProcedure framework is designed to support extensibility across all
three stages of the learning pipeline: K-Step (divergence-based clustering),
F-Step (cluster-wise learning), and C-Step (ensemble aggregation). Each stage
exposes a stable interface that enables substitution of algorithmic components
without modifying the core pipeline logic.

Extensibility is achieved through a factory and registry mechanism that
dynamically selects implementations based on configuration. This design
allows the framework to integrate new learning strategies in a modular and
decoupled manner.

\begin{figure}[H]
    \centering

    \resizebox{0.98\linewidth}{!}{%
    \begin{tikzpicture}[
        box/.style={
            draw,
            rounded corners,
            align=center,
            text width=4.3cm,
            minimum height=1.3cm,
            font=\small
        },
        mainbox/.style={box, fill=gray!10},
        newbox/.style={box, fill=gray!5},
        arrow/.style={-Latex, thick}
    ]

        % Core system
        \node[mainbox] (system) {KFCProcedure System\\
        \footnotesize K-Step $\rightarrow$ F-Step $\rightarrow$ C-Step};

        % Interface layer
        \node[box, below=1.2cm of system] (interface)
        {Abstract Interfaces\\
        \footnotesize Divergence / Estimator / Combiner};

        % Factory layer
        \node[box, below=1.2cm of interface] (factory)
        {Factory \& Registry Layer\\
        \footnotesize configuration-driven selection of components};

        % Existing components (KFC-based only)
        \node[box, below left=1.5cm and 1.2cm of factory] (existing)
        {Existing Components\\
        \footnotesize
        K-Step: Euclidean\\
        F-Step: linear \\
        C-Step: voting combiners};

        % New extensions (still KFC scope)
        \node[newbox, below right=1.5cm and 1.2cm of factory] (new)
        {Pluggable Extensions\\
        \footnotesize
        new divergences for K-Step\\
        new estimators for F-Step\\
        new combiners for C-Step};

        % Connections
        \draw[arrow] (system) -- (interface);
        \draw[arrow] (interface) -- (factory);
        \draw[arrow] (factory) -- (existing);
        \draw[arrow] (factory) -- (new);

    \end{tikzpicture}%
    }

    \caption{Extensibility in KFCProcedure across K-Step, F-Step, and C-Step components}
\end{figure}

\textbf{Explanation.}
The extensibility of the KFCProcedure framework is achieved through a
layered design in which each stage of the pipeline defines abstract
interfaces for its functionality. The K-Step supports interchangeable
divergence measures for clustering, the F-Step allows different supervised
learning models as local estimators, and the C-Step enables multiple
aggregation strategies for final prediction.

A factory and registry mechanism is used to instantiate concrete
implementations based on configuration parameters. This enables seamless
integration of new algorithms within each stage of the pipeline while
preserving the stability of the overall architecture.

As a result, the framework supports modular experimentation across all
components of the KFC pipeline without requiring modification of existing
code.

% ============================================================
\subsection{Summary}

Overall, the \texttt{kfc\_procedure} framework is built on modular,
interface-driven, and extensible design principles. The combination of
layered architecture, composition, and design patterns ensures that the
system remains flexible, scalable, and suitable for both research and
production environments.

\section{Core Interface Layer}
\label{sec:core_interface_layer}

The repository does not implement one single universal interface layer.
Instead, it defines several interface families for different parts of
the framework. The KFC subsystem defines abstractions for divergences,
local models, and combiners, while the COBRA subsystem defines
abstractions for estimators, splitters, distances, kernels, losses,
optimizers, normalizers, adapters, cross-validation, and aggregators.

This design keeps the two subsystems modular while still allowing them
to interact through the prediction matrix passed from the KFC F-step to
the C-step.

\subsection{Factory Infrastructure}

Two registry-based factory infrastructures are implemented:

\begin{itemize}
      \item \texttt{core/factory.py}: used by KFC components such as
            divergence factories, local model factories, and combiner factories.

      \item \texttt{cobra/core/factory.py}: used by COBRA components such
            as kernels, distances, losses, optimizers, splitters, estimators,
            adapters, normalizers, and aggregators.
\end{itemize}

Both factory systems support case-insensitive registration, aliases,
category labels, runtime creation, and discovery of available
components. This allows users to select components by string identifiers
such as \texttt{"euclidean"}, \texttt{"rbf"}, \texttt{"mse"},
\texttt{"grid"}, or \texttt{"weighted\_mean"}.

\subsection{KFC Interface Families}

The KFC subsystem relies on the following core abstractions:

\begin{itemize}
      \item \textbf{\texttt{BaseBregmanDivergence}:} Defines the interface
            for divergence functions used by \texttt{BregmanKMeans}. Concrete
            divergences implement \texttt{phi()}, \texttt{grad\_phi()}, and
            \texttt{in\_domain()}.

      \item \textbf{\texttt{BregmanDivergenceFactory}:} Registers and
            instantiates divergence implementations. The registered divergence
            names are \texttt{"euclidean"}, \texttt{"gkl"}, \texttt{"is"}, and
            \texttt{"logistic"}.

      \item \textbf{\texttt{BaseLocalModel}:} Defines the local model
            interface used by the F-step. Local models must implement
            \texttt{fit()} and \texttt{predict()}.

      \item \textbf{\texttt{LocalModelFactory}:} Registers local models.
            It includes a native mean regressor and automatically registers
            compatible scikit-learn classifiers and regressors.

      \item \textbf{\texttt{BaseCombiner}:} Defines the interface for
            C-step combiners. Each combiner implements \texttt{fit()} and
            \texttt{combine()}, while \texttt{predict()} delegates to
            \texttt{combine()}.

      \item \textbf{\texttt{CombinerFactory}:} Registers regression and
            classification combiners, including simple rules, stacking
            combiners, and COBRA-based combiners.
\end{itemize}

The clustering workflow itself is implemented by \texttt{BregmanKMeans}
and \texttt{KStep}. There is no separate generic \texttt{BaseClusterer}
class in the repository.

\subsection{COBRA Interface Families}

The COBRA subsystem defines a larger set of interfaces under
\texttt{cobra/core/}. These include:

\begin{itemize}
      \item \textbf{\texttt{BaseEstimator}} and
            \textbf{\texttt{EstimatorFactory}} for COBRA base learners.

      \item \textbf{\texttt{BaseDataSplitter}} and
            \textbf{\texttt{SplitterFactory}} for train/calibration splitting.

      \item \textbf{\texttt{BaseDistance}} and
            \textbf{\texttt{DistanceFactory}} for pairwise distance computation.

      \item \textbf{\texttt{BaseKernel}} and
            \textbf{\texttt{KernelFactory}} for distance-to-weight
            transformations.

      \item \textbf{\texttt{BaseKernelAdapter}} and
            \textbf{\texttt{KernelAdapterFactory}} for one-parameter and
            two-parameter distance transformations.

      \item \textbf{\texttt{BaseLoss}} and
            \textbf{\texttt{LossFactory}} for optimization objectives.

      \item \textbf{\texttt{BaseOptimizer}} and
            \textbf{\texttt{OptimizerFactory}} for grid-search and
            gradient-based parameter optimization.

      \item \textbf{\texttt{BaseAggregator}} and
            \textbf{\texttt{AggregatorFactory}} for weighted mean and weighted
            vote aggregation.

      \item \textbf{\texttt{BaseCrossValidator}} and
            \textbf{\texttt{CVFactory}} for K-fold, stratified K-fold, and
            time-series cross-validation strategies.

      \item \textbf{\texttt{BaseNormalizer}} and
            \textbf{\texttt{NormalizerFactory}} for reusable standard and
            min-max normalization utilities.
\end{itemize}

\subsection{Scikit-Learn Compatibility}

The high-level estimators inherit from or follow scikit-learn-style
interfaces. \texttt{KFCProcedure}, \texttt{KFCRegressor}, and
\texttt{KFCClassifier} expose \texttt{fit()} and \texttt{predict()}.
The COBRA classes \texttt{GradientCOBRA} and
\texttt{MixCOBRARegressor} inherit from scikit-learn estimator and
regressor mixin classes, while \texttt{CombinedClassifier} follows a
classifier-style interface with \texttt{fit()}, \texttt{predict()}, and
\texttt{predict\_proba()}.

One implementation limitation should be noted: the full
\texttt{KFCProcedure.predict\_proba()} method calls
\texttt{FStep.predict\_proba()}, but the current \texttt{FStep} class
does not implement this method. Therefore, full KFC probability
prediction should not be described as a completed feature in the current
implementation.

\section{Clusterwise Learning Subsystem (KFC)}
\label{sec:kfc_subsystem}

The Clusterwise Learning Subsystem implements the KFC part of the
library. It is located under \texttt{src/kfc\_procedure/core/} and is
responsible for divergence-based clustering, local model training, and
combiner-based prediction fusion.

\subsection{Subsystem Architecture}

The KFC subsystem is organized into five main packages:

\begin{verbatim}
core/
├── clustering/
├── ml/
├── steps/
├── combiner/
└── factory.py
\end{verbatim}

The implementation separates the clusterwise workflow into three
pipeline stages:

\begin{itemize}
      \item \textbf{K-step:} implemented in \texttt{steps/kstep.py};
            fits one \texttt{BregmanKMeans} model per divergence.

      \item \textbf{F-step:} implemented in \texttt{steps/fstep.py};
            trains local models for each divergence and cluster.

      \item \textbf{C-step:} implemented in \texttt{steps/cstep.py};
            combines the divergence-level prediction matrix using a registered
            combiner.
\end{itemize}

\subsection{Clustering Module}

The clustering module is implemented in:

\begin{verbatim}
core/clustering/
├── bregman.py
└── divergences/
\end{verbatim}

The central class is \texttt{BregmanKMeans}. It follows a Lloyd-style
clustering procedure with the following implementation features:

\begin{itemize}
      \item random centroid initialization from training samples;
      \item multiple random initializations through \texttt{n\_init};
      \item assignment using the selected Bregman divergence;
      \item centroid update using arithmetic means;
      \item empty-cluster reinitialization;
      \item distortion computation, including a streaming variant for
            memory-safe distance evaluation;
      \item scikit-learn-style methods such as \texttt{fit()},
            \texttt{predict()}, \texttt{transform()}, \texttt{fit\_predict()},
            \texttt{fit\_transform()}, and \texttt{score()}.
\end{itemize}

The divergence implementations are located in
\texttt{core/clustering/divergences/}. The registered divergences are:

\begin{itemize}
      \item \texttt{"euclidean"}: \texttt{SquaredEuclidean};
      \item \texttt{"gkl"}: \texttt{GKLDivergence};
      \item \texttt{"is"}: \texttt{ItakuraSaito};
      \item \texttt{"logistic"}: \texttt{LogisticLoss}.
\end{itemize}

The \texttt{validate\_divergence\_domain()} function checks that the
input matrix contains finite values and satisfies the selected
divergence domain before clustering is performed.

\subsection{Machine Learning Module}

The local learning module is implemented under \texttt{core/ml/}. It
contains:

\begin{itemize}
      \item \texttt{BaseLocalModel}: the abstract interface for F-step
            local models;
      \item \texttt{MeanRegressor}: a native local model based on
            scikit-learn's \texttt{DummyRegressor};
      \item \texttt{SklearnLocalModel}: an adapter that wraps compatible
            scikit-learn estimators;
      \item \texttt{LocalModelFactory}: a registry for resolving local
            models by name.
\end{itemize}

The function \texttt{register\_all\_sklearn\_models()} automatically
registers compatible scikit-learn classifiers and regressors using
snake-case names such as \texttt{linear\_regression},
\texttt{logistic\_regression}, \texttt{random\_forest\_regressor}, and
\texttt{k\_neighbors\_classifier}, depending on the classes available
from scikit-learn.

\subsection{KFC Workflow Steps}

The KFC pipeline is implemented through three concrete step classes:

\begin{itemize}
      \item \textbf{\texttt{KStep}:} Receives a list of divergence
            identifiers or divergence objects. It creates and fits one
            \texttt{BregmanKMeans} model per divergence and stores both fitted
            models and training cluster assignments.

      \item \textbf{\texttt{FStep}:} Receives the cluster labels generated
            by \texttt{KStep}. For each divergence and each observed cluster, it
            resolves a local model from \texttt{LocalModelFactory}, trains it on
            the corresponding subset, and stores it in a nested dictionary.

      \item \textbf{\texttt{CStep}:} Receives the prediction matrix
            generated by \texttt{FStep}. It resolves a combiner from
            \texttt{CombinerFactory}, fits it when required, and exposes
            \texttt{predict()} for final aggregation.
\end{itemize}

During prediction, \texttt{FStep.predict()} returns a matrix of shape
\texttt{(n\_samples, n\_divergences)}. Each column corresponds to one
divergence-specific predictive view. The clusters are used internally to
route each sample to the correct local model, but the final matrix has
one column per divergence, not one column per cluster.

\subsection{Combiner Module}

The combiner module is implemented under \texttt{core/combiner/}. It
contains separate regression and classification combiners.

The registered regression combiners are:

\begin{itemize}
      \item \texttt{"mean"}: row-wise arithmetic mean;
      \item \texttt{"weighted\_mean"}: linear-regression-based weighted
            combination;
      \item \texttt{"stacking\_regressor"}: meta-regression using a
            configurable regressor, with linear regression as the default;
      \item \texttt{"gradientcobra"}: wrapper around
            \texttt{cobra.gradientcobra.GradientCOBRA};
      \item \texttt{"mixcobra"}: wrapper around
            \texttt{cobra.mixcobra.MixCOBRARegressor}.
\end{itemize}

The registered classification combiners are:

\begin{itemize}
      \item \texttt{"majority\_vote"}: hard voting across prediction
            columns;
      \item \texttt{"stacking\_classifier"}: logistic-regression-based
            stacking;
      \item \texttt{"combined\_classifier"}: wrapper around
            \texttt{cobra.combined\_classifier.CombinedClassifier}.
\end{itemize}

The COBRA-based KFC combiners call the corresponding COBRA model with
\texttt{as\_predictions=True}. This is important because the input to
the combiner is already a prediction matrix produced by the F-step.

\subsection{Subsystem Summary}

The KFC subsystem implements the clusterwise learning part of the
framework through explicit packages for clustering, local model
wrapping, pipeline steps, and prediction combiners. Its implementation
is modular: divergence selection, local model choice, and final
aggregation strategy can be changed independently through the factory
registries.


\section{COBRA Aggregation Subsystem}
\label{sec:cobra_subsystem}

The COBRA Aggregation Subsystem implements prediction-space aggregation
models and reusable aggregation components. It is located under
\texttt{src/kfc\_procedure/cobra/}. Unlike the KFC subsystem, which
first partitions the input space into clusters, the COBRA subsystem
works primarily by constructing or receiving a prediction matrix and
then aggregating calibration targets according to distance-based kernel
weights.

\subsection{Subsystem Architecture}

The COBRA package contains four main model-level files:

\begin{itemize}
      \item \texttt{gradientcobra.py}: implements
            \texttt{GradientCOBRA} for regression aggregation.

      \item \texttt{mixcobra.py}: implements
            \texttt{MixCOBRARegressor}, which combines input-space and
            prediction-space distances.

      \item \texttt{combined\_classifier.py}: implements
            \texttt{CombinedClassifier} and \texttt{CombinedClassifierFast} for
            classification aggregation.

      \item \texttt{superlearner.py}: implements a standalone
            \texttt{SuperLearner} stacking-style estimator.
\end{itemize}

The reusable component infrastructure is located under
\texttt{cobra/core/}. It contains adapters, aggregators,
cross-validation strategies, distances, estimators, kernels, losses,
normalizers, optimizers, splitters, factories, and shared type
definitions.

\subsection{Core Engine Structure}

The COBRA core package is implemented as a factory-driven component
system:

\begin{itemize}
      \item \textbf{Adapters:} \texttt{OneParameterKernelAdapter}
            scales one distance matrix by a bandwidth parameter.
            \texttt{TwoParameterKernelAdapter} combines two distance matrices
            using coefficients $\alpha$ and $\beta$.

      \item \textbf{Aggregators:} \texttt{WeightedMeanAggregator}
            performs regression aggregation, while
            \texttt{WeightedVoteAggregator} performs classification aggregation.

      \item \textbf{Cross-validation:} \texttt{KFoldCV},
            \texttt{StratifiedKFoldCV}, and \texttt{TimeSeriesCV} provide fold
            generators. The main COBRA classes use K-fold cross-validation by
            default.

      \item \textbf{Distances:} Euclidean, Manhattan, Minkowski, Cosine,
            and Hamming distances are implemented. Hamming distance is especially
            relevant for classification prediction vectors.

      \item \textbf{Estimators:} The COBRA estimator layer includes
            \texttt{BaseEstimator}, \texttt{MeanRegressor},
            \texttt{SklearnEstimator}, and automatic scikit-learn estimator
            registration through \texttt{EstimatorFactory}.

      \item \textbf{Kernels:} The registered kernels include radial/RBF,
            exponential, reverse-cosh, Cauchy, Epanechnikov, biweight,
            triweight, triangular, COBRA threshold, and naive kernels.

      \item \textbf{Losses:} The available losses include MSE, MAE,
            Huber, quantile, log loss, and hinge loss.

      \item \textbf{Normalizers:} Standard and min-max normalizers are
            implemented as reusable components. The main COBRA classes also use
            scalar normalization constants from \texttt{cobra/utils/preprocessing.py}.

      \item \textbf{Optimizers:} The optimizer layer includes
            \texttt{GridSearchOptimizer}, \texttt{GradientDescentOptimizer},
            \texttt{MomentumOptimizer}, and \texttt{AdamOptimizer}. These
            optimizers tune aggregation parameters such as bandwidth, alpha, and
            beta.

      \item \textbf{Splitters:} \texttt{RandomHoldoutSplitter} and
            \texttt{OverlapSplitter} create training and calibration partitions.
\end{itemize}

\subsection{GradientCOBRA Implementation}

\texttt{GradientCOBRA} is implemented as a scikit-learn-compatible
regressor. Its default estimator pool includes linear regression, ridge
CV, lasso CV, k-nearest-neighbor regression, random forest regression,
and support vector regression.

The \texttt{fit()} method supports two execution modes:

\begin{itemize}
      \item \textbf{Raw-feature mode:} The model splits the data into
            estimator-training and calibration subsets, trains the base
            estimators, and constructs a prediction matrix on the calibration
            data.

      \item \textbf{Prediction-matrix mode:} When
            \texttt{as\_predictions=True}, the input is treated as an existing
            prediction matrix and base estimator training is skipped.
\end{itemize}

After the prediction matrix is prepared, the class computes a scalar
normalization constant, constructs a distance matrix, applies a
one-parameter kernel adapter, evaluates kernel weights, and optimizes
the bandwidth through cross-validation.

\subsection{MixCOBRARegressor Implementation}

\texttt{MixCOBRARegressor} extends the prediction-space aggregation
idea by using both input-space and prediction-space distances. Its
default estimator pool includes linear regression, ridge, lasso,
k-nearest-neighbor regression, random forest regression, and support
vector regression.

The class supports two modes controlled by \texttt{one\_parameter}:

\begin{itemize}
      \item \textbf{One-parameter mode:} Input features and prediction
            features are concatenated after normalization, and a single
            bandwidth-like parameter is optimized.

      \item \textbf{Two-parameter mode:} Input-space distance and
            prediction-space distance are computed separately and combined by
            the two-parameter adapter using $\alpha$ and $\beta$.
\end{itemize}

Therefore, MixCOBRA should be described as a mixed-space distance
aggregation method, not as a method that combines multiple kernel
strategies.

\subsection{CombinedClassifier Implementation}

\texttt{CombinedClassifier} implements a COBRA-style classifier. Its
default estimator pool includes logistic regression, decision tree
classifier, support vector classifier, and k-nearest-neighbor
classifier. The default distance is Hamming distance, the default kernel
is RBF, and the default aggregator is weighted vote.

The class supports both raw-feature mode and
\texttt{as\_predictions=True} mode. In raw-feature mode, it trains base
classifiers and constructs prediction vectors. In prediction-matrix
mode, the given input matrix is used directly as the classifier
prediction space.

The file also contains \texttt{CombinedClassifierFast}, an optional
extension that can build a FAISS index when FAISS is available and
\texttt{use\_faiss=True}. This is an acceleration path for nearest
neighbor search, not the default classifier path.

\subsection{SuperLearner Implementation}

The repository includes a standalone \texttt{SuperLearner} estimator in
\texttt{cobra/superlearner.py}. This class implements a stacking-style
workflow with base learners, cross-validated prediction features, and
meta-learners. It also supports \texttt{as\_predictions=True}, where the
input is treated directly as prediction features.

Unlike the newer COBRA core components, \texttt{SuperLearner} is less
integrated with the factory-based COBRA core infrastructure. It should
therefore be described as a standalone ensemble component rather than as
a core dependency of \texttt{GradientCOBRA} or \texttt{MixCOBRARegressor}.

\subsection{Subsystem Summary}

The COBRA subsystem implements a modular prediction-space aggregation
framework. Its main classes reuse core components for splitting,
estimator management, normalization, distance computation, kernel
weighting, optimization, loss evaluation, and aggregation. The same
subsystem can be used independently or as a KFC combiner through the
C-step.

\section{Orchestration Layer (kfc.py)}
\label{sec:orchestration_layer}

The orchestration layer is implemented in \texttt{kfc.py}. It provides
the high-level estimator classes \texttt{KFCProcedure},
\texttt{KFCRegressor}, and \texttt{KFCClassifier}. This layer coordinates
the K-step, F-step, and C-step into one end-to-end workflow.

\subsection{Main Classes}

The orchestration layer exposes three classes:

\begin{itemize}
      \item \texttt{KFCProcedure}: the general estimator with a
            \texttt{task} parameter set to either \texttt{"regression"} or
            \texttt{"classification"}.

      \item \texttt{KFCRegressor}: a convenience wrapper that initializes
            \texttt{KFCProcedure} with \texttt{task="regression"}.

      \item \texttt{KFCClassifier}: a convenience wrapper that initializes
            \texttt{KFCProcedure} with \texttt{task="classification"}.
\end{itemize}

\subsection{Training Execution Flow}

The \texttt{fit()} method performs the following verified sequence:

\begin{enumerate}
      \item Convert \texttt{X} and \texttt{y} into NumPy arrays.
      \item Split the dataset into two subsets using
            \texttt{train\_test\_split} with \texttt{test\_size=0.5}.
      \item Use stratified splitting when \texttt{task="classification"}.
      \item Fit \texttt{KStep} on the K/F training subset.
      \item Obtain training cluster assignments from
            \texttt{kstep\_.clusters\_}.
      \item Predict cluster assignments for the C-step calibration subset.
      \item Fit \texttt{FStep} using the K/F subset and training cluster
            assignments.
      \item Generate the calibration prediction matrix using
            \texttt{FStep.predict()}.
      \item Fit \texttt{CStep} on the calibration prediction matrix and
            its corresponding target values.
\end{enumerate}

This implementation uses the first split for clustering and local model
training, and the second split for training the final combiner. The
purpose is to avoid fitting the combiner on the same predictions used to
train the local models.

\subsection{Prediction Execution Flow}

The \texttt{predict()} method performs the following sequence:

\begin{enumerate}
      \item Check that \texttt{kstep\_}, \texttt{fstep\_}, and
            \texttt{cstep\_} are fitted.
      \item Convert the new input \texttt{X} into a NumPy array.
      \item Use \texttt{KStep.predict()} to assign clusters under each
            fitted divergence.
      \item Use \texttt{FStep.predict()} to route samples to local models
            and build the divergence-level prediction matrix.
      \item Use \texttt{CStep.predict()} to produce the final prediction.
\end{enumerate}

The orchestration layer does not perform general-purpose preprocessing
such as scaling, missing-value imputation, or encoding. These steps are
handled in notebooks or by the user before calling the framework.

\subsection{Combiner Integration}

The C-step combiner may be a simple rule-based combiner, a stacking
combiner, or a COBRA-based combiner. Therefore, the orchestration layer
should not be described as always ending with COBRA aggregation. COBRA
aggregation is used only when the selected combiner is
\texttt{"gradientcobra"}, \texttt{"mixcobra"}, or
\texttt{"combined\_classifier"}.

\subsection{Probability Prediction Limitation}

Although \texttt{KFCProcedure} defines a \texttt{predict\_proba()} method
for classification, the current implementation calls
\texttt{FStep.predict\_proba()}, which is not implemented in
\texttt{FStep}. Therefore, Chapter~6 should not present full KFC
probability prediction as a completed supported workflow. Probability
prediction is implemented inside some lower-level components, such as
\texttt{CombinedClassifier}, but not through the full KFC orchestration
pipeline in its current state.

\subsection{Subsystem Summary}

The orchestration layer provides a unified entry point for KFC training
and inference. It coordinates the three pipeline stages while allowing
the user to select divergences, local models, and combiners through
configuration parameters and factory-registered component names.

\section{Utility Layer}
\label{sec:utility_layer}

The repository contains two utility areas: the top-level
\texttt{kfc\_procedure/utils/} package and the COBRA-specific
\texttt{kfc\_procedure/cobra/utils/} package. These utilities support
pipeline execution but do not implement the main learning algorithms.

\subsection{Top-Level KFC Utilities}

The top-level utility package contains:

\begin{itemize}
      \item \texttt{logger.py}: defines a lightweight \texttt{Logger}
            class with verbosity levels for KFC execution messages. It provides
            \texttt{info()}, \texttt{debug()}, and \texttt{trace()} methods, as
            well as a \texttt{timed()} decorator.

      \item \texttt{resolve.py}: provides helper functions for resolving
            Bregman clustering configurations into \texttt{BregmanKMeans}
            objects. This module contains \texttt{resolve\_bregman()} and
            \texttt{resolve\_kstep()}.
\end{itemize}

The current top-level utility package does not contain a general
preprocessing pipeline for missing-value handling, scaling, or encoding.

\subsection{COBRA Utilities}

The COBRA utility package contains:

\begin{itemize}
      \item \texttt{preprocessing.py}: implements
            \texttt{data\_split\_overlap()}, \texttt{compute\_normalization\_constant()},
            \texttt{clean\_sklearn\_name()}, and
            \texttt{history\_to\_dataframe()}.

      \item \texttt{resolve.py}: implements helper functions for resolving
            estimators, kernels, splitters, losses, distances, aggregators, and
            training contexts. It also contains \texttt{fit\_estimators()} and
            \texttt{predict\_estimators()}.

      \item \texttt{distance.py}: implements a Numba-accelerated Hamming
            distance matrix function.
\end{itemize}

These utilities are used mainly by the COBRA model classes to construct
training/calibration contexts, fit estimator pools, compute prediction
matrices, normalize prediction spaces, and resolve factory-based
components.

\subsection{Subsystem Summary}

The utility layer provides lightweight support functions for logging,
configuration resolution, splitting, normalization constants, and
estimator management. It should be described as a support layer rather
than a full preprocessing or configuration-management subsystem.


\section{Training and Prediction Workflows}
\label{sec:training_prediction_workflows}

This section describes the implemented workflows for the two main usage
patterns in the repository: the KFCProcedure workflow and the direct
COBRA workflow.

\subsection{KFCProcedure Training Workflow}

The KFC training workflow is implemented in \texttt{KFCProcedure.fit()}.
The verified sequence is:

\begin{enumerate}
      \item Receive input features \texttt{X} and target vector \texttt{y}.
      \item Convert both inputs into NumPy arrays.
      \item Split the data into \texttt{(X\_k, y\_k)} and
            \texttt{(X\_l, y\_l)} using a 50/50 split.
      \item Fit \texttt{KStep} on \texttt{X\_k}.
      \item Use the fitted K-step to obtain cluster labels for both
            \texttt{X\_k} and \texttt{X\_l}.
      \item Fit \texttt{FStep} on \texttt{X\_k}, \texttt{y\_k}, and the
            K-step training cluster labels.
      \item Generate the calibration prediction matrix
            \texttt{P\_l = FStep.predict(X\_l, clusters\_l)}.
      \item Fit \texttt{CStep} on \texttt{P\_l} and \texttt{y\_l}.
\end{enumerate}

The C-step may use a simple combiner, a stacking combiner, or a
COBRA-based combiner, depending on the selected configuration.

\subsection{KFCProcedure Prediction Workflow}

The prediction workflow is implemented in \texttt{KFCProcedure.predict()}:

\begin{enumerate}
      \item Check that the K-step, F-step, and C-step have been fitted.
      \item Predict divergence-specific cluster assignments for the new
            input matrix.
      \item Route each sample through the corresponding cluster-local model
            in the F-step.
      \item Stack the outputs into a prediction matrix with one column per
            divergence.
      \item Pass the prediction matrix to the fitted C-step combiner.
      \item Return the final prediction vector.
\end{enumerate}

\subsection{Direct COBRA Training Workflow}

The COBRA classes can also be used directly without KFC. In raw-feature
mode, \texttt{GradientCOBRA}, \texttt{MixCOBRARegressor}, and
\texttt{CombinedClassifier} use the following general workflow:

\begin{enumerate}
      \item Resolve a training/calibration context.
      \item Fit a pool of base estimators on the training subset.
      \item Predict on the calibration subset to construct a prediction
            matrix.
      \item Compute distances in prediction space, or in both input and
            prediction spaces for MixCOBRA.
      \item Transform distances using a kernel adapter.
      \item Apply a kernel function to produce weights.
      \item Optimize aggregation parameters using cross-validation.
      \item Store calibration targets and fitted components for
            prediction.
\end{enumerate}

\subsection{Prediction-Matrix Mode}

The COBRA classes support \texttt{as\_predictions=True}. In this mode,
the input matrix is treated as an already computed prediction matrix.
Base estimator training is skipped, and the model directly performs
distance computation, kernel weighting, optimization, and aggregation on
the supplied prediction matrix.

This mode is used by the KFC COBRA-based combiners because the F-step
already generates a prediction matrix.

\subsection{Workflow Summary}

The framework supports both clusterwise learning and direct
prediction-space aggregation. KFC creates divergence-level predictions
through clustering and local models, while COBRA aggregates prediction
matrices using distance, kernel, optimizer, loss, and aggregator
components.


\section{Testing and Validation}
\label{sec:testing_validation}

The repository includes an automated test suite under \texttt{tests/}.
The verified tests currently focus on the COBRA core components. No
automated KFC subsystem tests were found in the repository at the time
of this audit.

\subsection{Test Suite Structure}

The test suite is organized as follows:

\begin{verbatim}
tests/
└── cobra/
    ├── test_p01_splitters.py
    ├── test_p02_estimators.py
    ├── test_p03_normalizers.py
    ├── test_p04_distances.py
    ├── test_p05_kernel_adapters.py
    ├── test_p06_kernels.py
    ├── test_p07_losses.py
    ├── test_p08_cv.py
    ├── test_p09_optimizers.py
    └── test_p10_aggregators.py
\end{verbatim}

\subsection{Covered Components}

The automated tests cover the following COBRA components:

\begin{itemize}
      \item \textbf{Splitters:} holdout splitting, overlap splitting,
            reproducibility, factory creation, and parameter validation.

      \item \textbf{Estimators:} mean regressor behavior, scikit-learn
            wrapper behavior, factory creation, and interface contract checks.

      \item \textbf{Normalizers:} standard normalization, min-max
            normalization, constant-feature handling, and factory resolution.

      \item \textbf{Distances:} Euclidean, Manhattan, Minkowski, Cosine,
            and Hamming distance behavior.

      \item \textbf{Kernel adapters:} one-parameter scaling,
            two-parameter distance combination, parameter updates, and input
            validation.

      \item \textbf{Kernels:} factory aliases and behavior of radial,
            exponential, reverse-cosh, Cauchy, Epanechnikov, biweight,
            triweight, triangular, naive, and COBRA kernels.

      \item \textbf{Losses:} MSE, MAE, Huber, log loss, hinge, and
            quantile loss.

      \item \textbf{Cross-validation:} K-fold fold counts, disjoint
            validation sets, full coverage, no train/validation leakage, and
            reproducibility.

      \item \textbf{Optimizers:} gradient descent, Momentum, Adam, grid
            search, optimizer factories, and optimization history.

      \item \textbf{Aggregators:} weighted mean, weighted vote, fallback
            behavior, NaN handling, batch aggregation, and empty-value checks.
\end{itemize}

\subsection{Current Testing Gaps}

The current repository does not contain separate automated tests for:

\begin{itemize}
      \item \texttt{BregmanKMeans} and divergence-specific clustering;
      \item \texttt{KStep}, \texttt{FStep}, and \texttt{CStep};
      \item full \texttt{KFCProcedure} training and prediction workflows;
      \item KFC regression and classification wrappers;
      \item notebook execution as an automated validation pipeline.
\end{itemize}

Therefore, the chapter should not claim complete system-level testing
unless additional tests are added later.

\subsection{Validation Through Notebooks}

The repository also includes experiment notebooks under
\texttt{notebooks/}. These notebooks demonstrate regression and
classification experiments using baseline models, COBRA models, and KFC
configurations. However, notebook execution is not currently integrated
into the automated pytest suite or CI workflow.

\subsection{Testing Summary}

The implemented automated tests provide good coverage for the COBRA
core component layer. The KFC subsystem and full orchestration workflow
are demonstrated in notebooks but are not yet covered by automated unit
or integration tests in the repository.

\section{Documentation and CI/CD Pipeline}
\label{sec:documentation_cicd}

This section describes the documentation and automation infrastructure
that is verified in the repository. The project includes README-level
documentation, inline docstrings, package metadata, pytest configuration,
and a GitHub Actions workflow for package testing, building, and PyPI
publishing.

\subsection{Documentation Assets}

The verified documentation assets are:

\begin{itemize}
      \item \texttt{README.md}: provides a project overview, feature list,
            installation instructions, quick-start examples, and high-level
            component descriptions.

      \item \textbf{Inline docstrings}: many modules and classes include
            explanatory docstrings describing parameters, workflows, and design
            intent.

      \item \texttt{pyproject.toml}: documents package metadata,
            dependencies, optional extras, setuptools configuration, and pytest
            settings.

      \item \texttt{notebooks/}: provides executable examples and
            experimental usage patterns for regression and classification tasks.
\end{itemize}

No separate Sphinx, MkDocs, ReadTheDocs, or GitHub Pages documentation
configuration was found in the repository. Therefore, automated
documentation generation should not be claimed unless such tooling is
added.

\subsection{Package Configuration}

The package is configured through \texttt{pyproject.toml}. The verified
configuration includes:

\begin{itemize}
      \item package name: \texttt{kfc-procedure};
      \item package source layout: \texttt{src/};
      \item Python requirement: \texttt{>=3.11};
      \item core dependencies such as NumPy, pandas, scikit-learn,
            matplotlib, and XGBoost;
      \item optional extras for \texttt{core}, \texttt{cobra},
            \texttt{dev}, and \texttt{all};
      \item pytest settings including \texttt{testpaths = ["tests"]}.
\end{itemize}

\subsection{Continuous Integration and Publishing Workflow}

The repository contains one GitHub Actions workflow:

\begin{verbatim}
.github/workflows/ci-cd.yml
\end{verbatim}

The workflow is named \texttt{Python Package CI/CD}. It is triggered by
Git tags matching the pattern:

\begin{verbatim}
v*.*.*
\end{verbatim}

The verified workflow steps are:

\begin{enumerate}
      \item check out the repository;
      \item set up Python 3.11;
      \item install dependencies using \texttt{pip install -e .[dev,cobra]};
      \item run the test suite using \texttt{pytest};
      \item build the package using \texttt{python -m build};
      \item upload the built distribution to PyPI using \texttt{twine} and
            a PyPI API token.
\end{enumerate}

This workflow is a tag-triggered package test/build/publish pipeline.
It is not configured to run on every push or pull request, and it does
not include verified linting, static analysis, or documentation
deployment steps.

\subsection{Automation Tools}

The verified automation tools are:

\begin{itemize}
      \item \textbf{GitHub Actions}: used for the tag-triggered workflow.
      \item \textbf{pytest}: used for automated tests.
      \item \textbf{setuptools}: used through the \texttt{pyproject.toml}
            package configuration.
      \item \textbf{build}: used to generate package distributions.
      \item \textbf{twine}: used to upload distributions to PyPI.
\end{itemize}

Although \texttt{.pylintrc} exists, the GitHub Actions workflow does not
currently run pylint, flake8, ruff, or another linter.

\subsection{Summary}

The repository includes a functional package configuration and a
tag-triggered CI/CD workflow for testing, building, and publishing the
package. Documentation is currently provided through the README,
docstrings, notebooks, and package metadata. More advanced documentation
deployment and broader CI checks would require additional workflow and
documentation configuration.

\section{Chapter Summary}
\label{sec:chapter_summary}

This chapter presented the implementation of the \texttt{kfc\_procedure}
framework as a modular Python research library. The repository combines
two related subsystems: the KFC clusterwise learning subsystem and the
COBRA prediction-space aggregation subsystem.

The KFC subsystem is implemented under \texttt{core/}. It includes
Bregman K-Means clustering, registered Bregman divergences, local model
wrappers, K-step/F-step/C-step workflow classes, and task-specific
combiner strategies. The COBRA subsystem is implemented under
\texttt{cobra/}. It includes \texttt{GradientCOBRA},
\texttt{MixCOBRARegressor}, \texttt{CombinedClassifier},
\texttt{SuperLearner}, and reusable core components for splitters,
estimators, distances, kernels, adapters, losses, optimizers,
normalizers, cross-validation, and aggregators.

The orchestration layer in \texttt{kfc.py} coordinates the KFC workflow
through \texttt{KFCProcedure}, \texttt{KFCRegressor}, and
\texttt{KFCClassifier}. During training, the implementation splits the
data into a K/F training subset and a C-step calibration subset, fits
divergence-specific clustering models, trains local models per cluster,
constructs a divergence-level prediction matrix, and fits the selected
combiner. During prediction, new samples are assigned to clusters,
routed through local models, and aggregated by the fitted C-step.

The repository applies modular software engineering principles through
factory registries, component-level abstractions, scikit-learn-style
interfaces, and clear separation between clustering, local modeling,
aggregation, optimization, and utility functions. The current automated
test suite focuses on COBRA core components, while KFC subsystem and
full workflow tests remain areas for future extension.

Finally, the project includes package metadata, README documentation,
experiment notebooks, and a GitHub Actions workflow that runs tests,
builds the package, and publishes to PyPI when a version tag is pushed.
Overall, the implementation translates the methodology from
Chapter~\ref{ch:methodology} into a reusable and extensible Python
framework while leaving room for additional automated testing and
documentation improvements.
